% =====================================================================
%  ONE-PAGE PAPER  —  style proposal
%  "Simulating Just-in-Time Adaptive Interventions: A Config-Driven,
%   LLM-Bootstrapped Framework"
%
%  One consistent format: 11pt everywhere, no size changes.
%  One-page deliverable: the report body fits on page 1; the reference list
%  follows on page 2 (deliberate scope, matching the compilation note below).
%  Compile: tectonic one_page_paper.tex
% =====================================================================
\documentclass[11pt,a4paper]{article}

% --- clean, tight, sleek layout ---
\usepackage[margin=1.5cm]{geometry}
\usepackage[T1]{fontenc}
\usepackage{microtype}
\usepackage{parskip}
\setlength{\parskip}{3pt}
\usepackage{booktabs}
\usepackage[dvipsnames]{xcolor}
\usepackage[numbers,sort&compress]{natbib}
\usepackage{titlesec}

% compact section titles
\titleformat{\section}{\bfseries}{\thesection.}{0.5em}{}
\titlespacing*{\section}{0pt}{6pt plus 2pt}{3pt plus 1pt}

% header rule under title
\newcommand{\paperrule}{\noindent\rule{\textwidth}{0.6pt}\par\vspace{6pt}}

% pipeline figure removed per style revision (2026-08-02)

\begin{document}

% ============================ PAGE 1 ================================
\begin{center}
  {\LARGE\bfseries Simulating Just-in-Time Adaptive Interventions:\\[2pt]
   A Config-Driven, LLM-Bootstrapped Framework}\\[4pt]
  {\normalsize William Dennis}\\[2pt]
  {\small University of Bristol \& National University of Singapore}\\[4pt]
  \paperrule
\end{center}

Physical inactivity is a leading modifiable health risk
\cite{guthold2018insufficient}; just-in-time adaptive
interventions (JITAIs) deliver personalised prompts through wearables, but trials
cannot probe counterfactuals and transition data is rarely public
\cite{wang2021optimizing}. We
reduce policy design to a config-driven pipeline for two simulators:
LLM-bootstrapped patient models, evaluated against an optimal dynamic-programming
bound and bandit baselines.

\section{Motivation and approach}
We formalise the intervention as a Markov decision process
$\mathcal{M}=(\mathcal{S},\mathcal{A},\mathcal{P},\mathcal{R})$ and build a
framework in which \emph{everything} is configuration: state spaces, action sets,
transition models, rewards, and agents are declared in YAML, so experiments differ by
config file alone (configs and tables in the repository)
\cite{klasnja2019heartsteps,trella2022designing}.

\section{LLM-bootstrapped transitions}
The transition function $P(s' \mid s, a)$ is the hardest component of any
patient model: it must encode how context, burden, and a delivered suggestion
change behaviour, and it is rarely available as data. Rather than hand-curating
tables or waiting for trials, we prompt an LLM (DeepSeek V4 Flash, with a GLM
5.2 comparison) to sample
next-state distributions given a persona description, current state, and
candidate action --- a text world model \cite{liwang2026textworld}. Unlike
data-driven simulators \cite{wang2021optimizing}, which require logged trial
data, bootstrapping needs
no logged data --- though it still requires hand-authored personas, states, and
rewards in config. Personas were hand-curated by domain practitioners, loosely
grounded in the behaviour-change literature: not arbitrary, but not validated
against real patient data either. \citet{park2024thousandpeople} report agents reproducing held-out
General Social Survey responses at 85\% of participants' own two-week response-replication
consistency; we make no such population-replication claim for ours. Six sweeps of 22{,}320 prompts each (133{,}920 prompts, 10 samples
per outcome) produced 31 JSON tables: \$24.52 (DeepSeek) + \$25.45 (GLM 5.2).
Personas produced distinct sampled transition distributions, and the config's
Bayesian mechanism models burden as lowering future activity probability
\cite{park2023generative,wang2021optimizing}.

\section{Validation}
On the Phase-1 persona tables, a HeartSteps-V2-style UCB (upper confidence
bound) bandit reaches 47--95.8\%
of the optimal dynamic-programming bound across personas (50 seeds; 95\%
confidence intervals $\pm$0.6--9.5pp);
the idle policy alone reaches 96.8--100.1\%, the values above 100\% being
Monte Carlo noise ($\pm$3--5) --- the bound is essentially the idle policy ---
interventions underperform idle on these tables, a property of these personas'
bootstrapped tables, not a claim about JITAI efficacy in general.
All baselines are internal to the bootstrapped tables (mock-on-mock;
no ground truth; burden sensitivity untested). Eight literature-review reports
document this background; the one quantitative comparison (a V2-style agent at
46.7\% vs the paper's 78.4\%, on different data) did not reproduce.

\section{PEARL-matched extension (Phase 2)}
In the final two weeks we compared against the PEARL randomised
controlled trial (RCT) \cite{lee2025pearl}, whose reinforcement-learning (RL)
arm gained $+$295.9 steps
over control (p=0.0002): we used
6 COM-B themes
$\times$ 2 delivery times \cite{michie2011com}, plus idle
--- with an $\varepsilon$-greedy bandit and a COM-B barrier-weighted fixed arm
\cite{liu2025thompson,ghosh2024rebandit,tomkins2021intelligentpooling}. The
14{,}040-prompt bootstrap (108 states $\times$ 13 actions $\times$ 10 samples)
passed 11 of 17 checks; the six failures
(inflated effect size $d{=}1{,}525$; fixed $>$ RL inversion; no attenuation;
no burden saturation; inverted weekend; degenerate control t-test) trace to
an idle-vs-action structure. A reward-penalty probe
($R=\mathrm{step\_reward}-\lambda\mathbf{1}[\mathrm{action}\ne\mathrm{idle}]$,
$\lambda=0.05$, step\_reward $-1/0/+1$; an anti-spam term from the Phase-1
reward design, independent of the PEARL comparison) does not restore the
RL-vs-fixed ordering: 50-seed runs sit within noise ($\pm$5--7) and the
ordering flips with the penalty (unpenalised RL $3.0>$ fixed $2.4$; penalised
fixed $-0.6>$ RL $-0.7$) --- the penalty taxes activeness. Calibration
metrics, not trial outcomes; the probe is not a validated fix.

\section{Conclusion}
\emph{Policy design reduces to a cheap, config-driven pipeline, albeit
mock-on-mock only, with a documented map of where the simulation diverges from
published trials.}
% ============================ PAGE 2+ ===============================
\clearpage
\begingroup
\bibliographystyle{unsrtnat}
\bibliography{../design/references}
\endgroup

\end{document}
